\documentclass[11pt,letterpaper]{article}
\usepackage[utf8]{inputenc}
\usepackage[T1]{fontenc}
\usepackage[spanish,es-tabla]{babel}
\usepackage[margin=2.4cm]{geometry}
\usepackage{amsmath,amssymb}
\usepackage{booktabs,longtable,array,tabularx}
\usepackage{enumitem}
\usepackage{xcolor}
\usepackage{hyperref}
\usepackage{fancyhdr}
\usepackage{microtype}
\hypersetup{colorlinks=true,linkcolor=blue!45!black,urlcolor=blue!45!black,citecolor=blue!45!black}
\setlength{\headheight}{15pt}
\pagestyle{fancy}
\fancyhf{}
\lhead{CubeSat UD34 --- Práctica 4}
\rhead{Central y comunicaciones}
\cfoot{\thepage}
\setlength{\parindent}{0pt}
\setlength{\parskip}{6pt}
\setlist[itemize]{noitemsep,topsep=3pt}
\renewcommand{\arraystretch}{1.18}
\newcolumntype{L}[1]{>{\raggedright\arraybackslash}p{#1}}

\begin{document}
\begin{titlepage}
\centering
{\large Universidad de Antioquia\par}
{Departamento de Ingeniería Electrónica y Telecomunicaciones\par}
\vspace{1.5cm}
{\Large Laboratorio de Diseño de Productos Electrónicos\par}
\vspace{0.6cm}
{\huge\bfseries Práctica 4\par}
\vspace{0.4cm}
{\LARGE Diseño electrónico y captura esquemática\par}
\vspace{0.8cm}
{\Large PCB integrada central y de comunicaciones del CubeSat UD34\par}
\vfill
\textbf{Integrantes}\par
Daniel Felipe Yépez Taimal\par
Santiago Cáceres Vásquez\par
\vspace{0.5cm}
\textbf{Docente}\par
Maria Katherine Plazas Olaya\par
\vspace{0.6cm}
Medellín, Colombia\par
23 de octubre de 2026\par
\vspace{0.5cm}

\end{titlepage}

\tableofcontents
\newpage
\section{Introducción}
El CubeSat UD34 se plantea como un sistema de formato 1U para adquirir imágenes multiespectrales de regiones cafeteras y transmitir información hacia una estación terrestre en Medellín. La práctica 3 estableció una arquitectura centralizada: un computador de a bordo coordina adquisición, almacenamiento, supervisión, comunicaciones y funciones de control de actitud \cite{lab3}.

La práctica 4 transforma esa arquitectura preliminar en un diseño electrónico detallado de la PCB asignada. Comprende especificaciones, componentes, dimensionamiento de circuitos, librerías, esquemáticos jerárquicos, restricciones, revisión eléctrica y costos. No exige implementar software de vuelo, enrutar la PCB ni generar archivos de fabricación \cite{guia4}.

La PCB integra el computador de a bordo, el almacenamiento y el subsistema de comunicaciones. El informe documenta sus funciones, sus interfaces internas y externas y las condiciones que sustentan el diseño electrónico. La integración requiere un presupuesto conjunto de alimentación, espacio y datos, así como compatibilidad eléctrica entre procesamiento y radiofrecuencia.

\section{Objetivos}
\subsection{Objetivo general}
Diseñar la PCB integrada central y de comunicaciones del CubeSat UD34, justificando la selección y el dimensionamiento de sus circuitos y verificando su correspondencia con los requisitos, las interfaces y las restricciones del sistema.
\subsection{Objetivos específicos}
\begin{itemize}
\item Definir las funciones de procesamiento, almacenamiento y comunicación de la PCB y sus conexiones con los demás subsistemas.
\item Revisar las especificaciones heredadas, justificar sus modificaciones y mantener la trazabilidad de las decisiones de diseño.
\item Establecer las necesidades de interfaces, procesamiento, almacenamiento, comunicación, alimentación y espacio.
\item Definir las interfaces eléctricas y la asignación de contactos de los conectores de la tarjeta.
\item Seleccionar y dimensionar los componentes y circuitos de la tarjeta con fundamento en sus condiciones de operación.
\item Capturar y verificar el esquemático jerárquico en Altium mediante símbolos y huellas validados.
\item Documentar las restricciones eléctricas y de fabricación, y estimar el costo de componentes para los lotes establecidos en la guía.
\item Registrar los cambios respecto de la práctica 3 y sus efectos sobre el sistema.
\end{itemize}

\section{Alcance y especificaciones de la tarjeta}
\label{sec:punto1}

\subsection{Alcance de la PCB integrada}
La PCB comprende el computador de a bordo, sus memorias de trabajo y programa, el almacenamiento de misión, el transceptor, los circuitos de radiofrecuencia y los elementos auxiliares de alimentación, control y supervisión. Sus funciones son:
\begin{enumerate}[label=\arabic*.]
\item Coordinar los modos de espera, adquisición, almacenamiento, transmisión y recuperación mediante recursos de hardware disponibles para software futuro.
\item Configurar la cámara y recibir sus datos a través de una interfaz compatible con el dispositivo seleccionado.
\item Conservar imágenes, telemetría y eventos en almacenamiento no volátil.
\item Gestionar la comunicación bidireccional con tierra mediante el transceptor y los circuitos de RF, intercambiando datos, comandos y estados con el computador de a bordo.
\item Intercambiar telemetría y señales de control con la tarjeta de potencia.
\item Proporcionar interfaces para sensores y controladores del sistema de actitud, cuyas etapas de potencia se consideran externas a esta PCB.
\item Permitir programación y depuración del computador, así como reinicio y comprobación eléctrica de los bloques de la PCB.
\end{enumerate}

Los límites externos de la PCB corresponden a la tarjeta de potencia, la cámara, los dispositivos de actitud y la antena. El cargador de baterías, los paneles solares, el diseño interno de la cámara y las etapas de potencia de actuadores pertenecen a esos subsistemas externos. La PCB incluye la conexión hacia antena y los circuitos de RF que requiera la solución de comunicaciones. La ubicación de sensores locales y circuitos de supervisión se determina según las necesidades de medición y las interfaces disponibles.

\subsection{Arquitectura funcional}
Se mantiene un computador principal como coordinador. El almacenamiento y la cámara se conectan a él; comunicaciones se conecta mediante un bus digital interno y señales auxiliares. Potencia y actitud permanecen como bloques externos cuyas interfaces eléctricas deben especificarse.

\begin{table}[htbp]
\centering
\caption{Ubicación de los bloques y relaciones funcionales principales.}
\label{tab:arquitectura}
\begin{tabularx}{\textwidth}{L{3.6cm}L{2.8cm}X}
\toprule Bloque & Ubicación & Relación funcional \\
\midrule
Computador de a bordo & PCB integrada & Coordina los bloques internos y los dispositivos externos. \\
Almacenamiento de misión & PCB integrada & Conserva los datos recibidos y los entrega al computador para su transmisión. \\
Comunicaciones y RF & PCB integrada & Intercambia datos y comandos con el computador y conecta con la antena. \\
Cámara y óptica & Externa a la PCB & Reciben configuración y entregan datos de imagen al computador. \\
Potencia & Tarjeta externa & Suministra energía y proporciona telemetría y alertas. \\
Actitud & Subsistema externo & Entrega mediciones y recibe comandos hacia sus controladores. \\
Antena & Externa a la PCB & Interfaz electromagnética del enlace con tierra. \\
\bottomrule
\end{tabularx}
\end{table}
La tabla distingue la ubicación física de los bloques de sus relaciones
funcionales. Los protocolos, niveles eléctricos y conectores se definirán
según los dispositivos seleccionados.

\subsection{Trazabilidad de requisitos aplicables}
La tabla conserva los identificadores de los requisitos de la práctica 3 \cite{lab3}. Su contenido expresa la contribución de la PCB a cada
requisito; no sustituye el enunciado original ni acredita su cumplimiento.
Las excepciones o modificaciones se documentarán explícitamente.
\small
\begin{longtable}{L{1.7cm}L{5.6cm}L{7.1cm}}
\toprule ID original & Aplicación a la PCB integrada & Evidencia prevista y frontera \\
\midrule\endhead
RF-01 & Adquirir y almacenar datos multiespectrales. & Interfaz de cámara compatible, recursos para transferencia y almacenamiento dimensionado. La cámara produce la imagen. \\
RF-02 & Transmitir datos y recibir comandos de tierra. & Interfaz digital interna compatible y enlace bidireccional dimensionado para el volumen de datos y los contactos previstos. \\
RF-03 & Recibir medidas y enviar comandos de actitud. & Interfaces hacia sensores y drivers externos; no se presupone la inclusión de actuadores en esta sección. \\
RF-04 & Acceder a telemetría eléctrica y alertas. & Interfaz con potencia y capacidad de recibir medidas. La medición física de batería y sus protecciones permanecen en potencia. \\
RNF-01 & Respetar el volumen del sistema 1U. & Reserva mecánica de la PCB integrada y de sus componentes; falta confirmar alturas y fijaciones. \\
RNF-02 & Mantener separación física de tarjetas. & PCB central/comunicaciones y PCB de potencia independientes; integrar central y radio conserva esa separación. \\
RNF-03 & Considerar integridad de las señales. & Restricciones según las interfaces realmente elegidas; no se exigirán pares diferenciales a buses que no los utilicen. \\
RA-01 & Considerar el ambiente de la misión. & Selección y documentación de condiciones operativas; el grado industrial no demuestra tolerancia a radiación o vacío. \\
RA-02 & Contribuir a la operación durante ciclos de iluminación y eclipse. & El presupuesto de consumo y la supervisión apoyan la gestión energética; la carga, descarga y protección de la batería corresponden a potencia. \\
RA-03 & Considerar integración y vibraciones. & Encapsulados, conectores y fijaciones apropiados al prototipo; sin afirmar calificación de lanzamiento. \\
RN-01 & Definir frecuencia y potencia compatibles con la estación terrestre y con los requisitos normativos aplicables. & Documentación del enlace y de las condiciones de operación. La banda y la potencia permanecen pendientes de verificación. \\
RN-02 & Incorporar recomendaciones de implementación. & Desacoplamiento, retornos, filtrado y restricciones extraídas posteriormente de fuentes del fabricante. \\
\bottomrule
\end{longtable}
\normalsize

\subsection{Revisión de especificaciones heredadas}
Las especificaciones ET-01 a ET-08 proceden de la práctica 3 \cite{lab3}. Se distingue su valor de referencia de las comprobaciones necesarias para adoptarlo en el diseño detallado.
\small
\begin{longtable}{L{1.4cm}L{4.1cm}L{8.9cm}}
\toprule ID & Base de la práctica 3 & Tratamiento en el diseño y justificación \\
\midrule\endhead
ET-01 & Batería 2S: 7,4 V nominal; 6,0--8,4 V. & Dato de contexto de potencia. No se conectará directamente a la lógica por asumir este rango; falta definir qué tensiones reguladas entrega la tarjeta de potencia. \\
ET-02 & Distribución de 3,3 V y 5 V, con tolerancia de $\pm5\,\%$. & Conservar como base preliminar y verificar los extremos 3,135--3,465 V y 4,75--5,25 V frente a los dispositivos seleccionados. Alimentación y nivel lógico se comprobarán por separado. \\
ET-03 & PCB de hasta $90\times90$ mm. & Aplicar a la PCB integrada completa. No asignar ese espacio por separado a central y comunicaciones. Confirmar con el equipo las zonas de conectores, fijaciones y altura. \\
ET-04 & Almacenamiento no volátil $\geq16$ GB. & Se adopta microSD de 16 GB nominales; exigir al menos 15 GB útiles tras formato. La retención de 90 días requiere 12,202 GB bajo las cuotas definidas. \\
ET-05 & Enlace preliminar 435--438 MHz. & Se conserva como antecedente para revisión en comunicaciones; no constituye por sí solo autorización de operación ni selección definitiva. \\
ET-06 & Tasa base $\geq9600$ bps. & Es un límite inferior preliminar, no una demostración de cumplimiento de la descarga diaria. Debe coordinarse con el volumen de datos y los contactos. \\
ET-07 & Operación entre $-40$ y $+85\,^{\circ}\mathrm{C}$. & Revisar contra condiciones de operación recomendadas, no máximos absolutos. Comprobar el rango en cada dispositivo de la PCB, incluido el almacenamiento. La cámara externa se evaluará por separado; cualquier relajación del requisito deberá quedar justificada. \\
ET-08 & Reloj del procesador $\geq80$ MHz. & Mantener trazabilidad, pero complementar con periféricos, RAM, DMA, pines y capacidad de transferencia. La frecuencia sola no demuestra desempeño. \\
\bottomrule
\end{longtable}
\normalsize

\subsection{Interfaces internas y externas de la PCB}
La siguiente tabla conserva los identificadores de la práctica 3 \cite{lab3}. Las alternativas de bus no son decisiones intercambiables ya resueltas: deben cerrarse usando las características del dispositivo de cada extremo.
\small
\begin{longtable}{L{2cm}L{3.1cm}L{5.5cm}L{3.3cm}}
\toprule ID & Relación & Definición preliminar & Condición de cierre \\
\midrule\endhead
IF-01, IF-05 & Potencia hacia central y memoria & Entradas reguladas de 5 V y 3,3 V y retorno común. & Corriente total, caída y secuencia; 3 A exclusivos del SoM. \\
IF-03 & Potencia hacia radio & Alimentación del bloque de comunicaciones dentro de la misma PCB. & Consumo de transmisión y asignación de redes. \\
IF-07 & Central hacia cámara & Configuración y captura por USB3, puerto USB\_2 anfitrión. & SDK y contrato de VBUS externo. \\
IF-08 & Cámara hacia central & Imágenes RAW16 por el mismo enlace USB3. & Tasa efectiva, buffers y formato. \\
IF-09 & Central y almacenamiento & Bus SD nativo de cuatro bits hacia microSD Kingston de misión; modo inicial 25 MHz y 3,3 V. & Circuito de alimentación, zócalo y escritura sostenida por verificar. \\
IF-10, IF-11 & Central y comunicaciones & SPI\_1 interno, Verdin maestro; MOSI, MISO, SCK y CS. & Referencia del extremo, niveles, modo y frecuencia. \\
IF-14, IF-15 & Central y actitud & UART\_1, 115200 bit/s, 8N1; controlador externo de actitud. & Tramas, tiempos y adaptación al dominio externo de 3,3 V. \\
IF-16 & Potencia y central & I2C\_1 a 100 kHz; EPS\_PGOOD y EPS\_ALERT\_N; dominio externo de 3,3 V. & Dirección, registros, pull-ups y estados sin alimentación. \\
IF-17, IF-18 & Supervisión y central & Medición/configuración de diagnóstico local o externo. & Ubicación y selección de sensores. \\
IF-12, IF-13 & Comunicaciones y estación terrestre & Enlace bidireccional a través de la cadena RF y la antena. & Banda, potencia, protocolo y capacidad útil. \\
Nueva: IF-C01 & Herramienta y central & USB\_1 para recuperación y UART\_3 para consola. & Conectores, niveles y estados de recuperación. \\
\bottomrule
\end{longtable}
\normalsize
Las interfaces IF-02 y IF-04 pertenecen a la alimentación de cámara y actitud, respectivamente, y la IF-06 a la supervisión. La alimentación de cámara y actuadores se resuelve externamente; se define su contrato de señales con la PCB integrada. La ubicación de la supervisión IF-06 se precisa según cada señal. IF-12 e IF-13 describen el enlace inalámbrico; su conexión física con la tarjeta se realiza mediante la interfaz hacia antena. Las funciones de control y transferencia de imágenes de IF-07 e IF-08 pueden compartir una misma interfaz física si la cámara seleccionada lo permite.

\subsection{Definición preliminar de conectores y señales}
\label{sec:conectores}
Se agrupan las señales por función. Se seleccionan conectores JST GH para actitud y telemetría de potencia, con contactos definidos en la sección de asignación de buses y conectores. La entrada de potencia se mantiene separada y requiere dimensionamiento de corriente. Los conectores USB, mantenimiento y antena se seleccionarán según sus requisitos eléctricos y de montaje.

\begin{tabularx}{\textwidth}{L{2.2cm}XX}
\toprule Identificación propuesta & Señales que se deben asignar & Criterio \\
\midrule
J\_EPS\_PWR & +5V\_IN, +3V3\_IN y retornos de potencia. & Conector y contactos por dimensionar con la corriente total por riel. \\
J\_EPS\_SIG & GND, EPS\_SCL, EPS\_SDA, EPS\_PGOOD y EPS\_ALERT\_N. & JST GH BM05B-GHS-TBT de cinco contactos; señales externas de 3,3 V. \\
J\_CAM & Pares de datos USB3, retorno y tratamiento explícito de VBUS. & Cámara alimentada externamente; topología de VBUS por cerrar sin retroalimentación. \\
J\_ADCS & GND, ADCS\_TX y ADCS\_RX, dominio externo de 3,3 V. & JST GH BM03B-GHS-TBT de tres contactos; potencia de actuadores externa. \\
J\_DBG & Acceso USB\_1 y consola UART\_3 con sus referencias. & Respetar niveles reales; definir si se usan conectores separados. \\
J\_ANT & Conexión RF entre la cadena de comunicaciones y la antena. & Tipo de conector, impedancia y límites de potencia según la solución RF. \\
Sin conector central--radio & Bus digital y señales auxiliares en redes internas. & Ambos bloques se encuentran en la misma PCB; evita introducir una separación física innecesaria. \\
\bottomrule
\end{tabularx}

\subsection{Contrato funcional de interfaces entre módulos}
\label{sec:contrato_interfaces}
La PCB integra computador, almacenamiento y comunicaciones. Potencia,
cámara y actitud se consideran módulos externos que satisfacen los
contratos aquí definidos. Esta partición permite desarrollar la
tarjeta sin diseñar sus circuitos externos, pero exige respetar las
tensiones, secuencias y condiciones de conexión en nuestra frontera.

\subsubsection{Alimentación y supervisión desde potencia}
Se requieren entradas reguladas de 5 V y 3,3 V, con retorno común.
La capacidad de 3 A establecida para el SoM pertenece únicamente a
su rama. La capacidad total de entrada debe incluir memoria externa,
comunicaciones, auxiliares y pérdidas. La tensión debe mantenerse
dentro de los límites de cada carga después de la distribución.

Se reserva \texttt{I2C\_1} para telemetría y comandos de potencia,
inicialmente a 100 kHz. El computador actúa como controlador y el
módulo de potencia como destino direccionado. Esta interfaz permite
consultar registros sin ocupar el SPI reservado a comunicaciones.
Se definen \texttt{EPS\_SCL}, \texttt{EPS\_SDA} y GND, además de
\texttt{EPS\_PGOOD} para supervisión de alimentación y
\texttt{EPS\_ALERT\_N} para avisos. La dirección I2C, los registros,
la temporización y el estado de esas señales sin alimentación se
establecerán antes de cerrar el esquemático. PGOOD debe actuar sobre
la supervisión apropiada; no se presupone que un programa pueda
supervisar una alimentación antes de que el computador arranque.

\subsubsection{Contrato con el módulo de actitud}
Se adopta un módulo externo con controlador propio que reúne las
mediciones y gobierna los drivers de los actuadores. El computador
intercambia comandos, consignas y estados con ese controlador; no se
conectan sensores sin interfaz directamente a una UART ni se traslada
la potencia de actuadores a nuestra PCB.

Se reserva \texttt{UART\_1}, inicialmente a 115200 bit/s, ocho bits
de datos, sin paridad y un bit de parada (8N1). Las señales se denominan
\texttt{ADCS\_TX} y \texttt{ADCS\_RX} desde la perspectiva del
computador: su TX se conecta al RX del módulo externo y viceversa.
La conexión punto a punto no necesita dirección de dispositivo.
Con un bit de inicio, ocho de datos y uno de parada, el techo de
carga útil antes de cabeceras es
\begin{equation}
R_{\mathrm{UART,max}}=\frac{115200}{1+8+1}
=11520\ \mathrm{bytes/s}.
\end{equation}
Este cálculo no demuestra suficiencia para un volumen de telemetría
todavía no definido. La longitud, identificación, secuencia y
comprobación de integridad de las tramas se especificarán junto con
los tiempos de respuesta. Se supone cableado corto dentro del 1U,
masa común y separación de retornos de actuadores; si esas condiciones
no permiten mantener integridad, se revisará la capa física.

\subsubsection{Niveles lógicos y adaptación}
Se establece un dominio de 3,3 V en las interfaces digitales externas
de potencia y actitud. Es un contrato de diseño de estos módulos,
no una propiedad supuesta de todos los pines del Verdin. Las señales
del SoM que operan a 1,8 V requieren adaptación compatible
\cite{verdinds}. Para I2C se debe conservar el funcionamiento de
drenador abierto y dimensionar las resistencias de polarización con
la capacitancia del bus y los límites eléctricos de ambos extremos.
Para UART se respeta la dirección de cada señal. Las señales de
supervisión también necesitan comprobar sus umbrales y estados.

Este acuerdo no modifica los niveles físicos de USB ni asigna
automáticamente 3,3 V al extremo interno del SPI de comunicaciones.
La referencia de alimentación de 5 V del SoM tampoco autoriza señales
lógicas de 5 V. La UART de actitud utiliza el TXU0202, desarrollado en la sección de
adaptación. Los adaptadores de I2C y, si procede, de SPI se seleccionarán
al cerrar sus condiciones eléctricas.

\subsubsection{SPI interno de comunicaciones}
Se reserva \texttt{SPI\_1}, con el Verdin como maestro, usando
\texttt{SPI\_MOSI}, \texttt{SPI\_MISO}, \texttt{SPI\_SCK} y
\texttt{SPI\_CS}. MOSI, SCK y CS salen del computador; MISO entra.
Se propone CS activo en bajo, que deberá confirmarse con el dispositivo
de comunicaciones. No se requiere un conector entre las dos secciones
de la misma PCB. La referencia de masa y los dominios de alimentación
deben quedar definidos aunque no se representen junto al bus en TOP.

La frecuencia, CPOL/CPHA, formato de transferencias y niveles del
extremo de comunicaciones se cerrarán con su referencia concreta.
\texttt{COM\_IRQ} y \texttt{COM\_RESET\_N} se reservan
conceptualmente; su necesidad y polaridad no se declaran confirmadas.
Los puertos de las hojas hijas deben coincidir con las entradas de
sus símbolos en TOP y tener direcciones coherentes. La continuidad
gráfica no sustituye la revisión eléctrica ni el ERC.

\subsubsection{Cámara alimentada externamente}
Se conserva \texttt{USB\_2} como anfitrión para la XIMEA. Control
y transferencia de imágenes comparten esa conexión. La entrega de
energía a la cámara se asigna al conjunto externo de potencia/cámara,
por lo que su consumo no se suma por defecto a la entrada de nuestra
PCB, aunque sí pertenece al presupuesto global del sistema.

La cámara seleccionada recibe energía mediante USB
\cite{ximeamanual}; asignar su alimentación externamente no crea una
entrada independiente en el dispositivo. El conjunto externo debe
resolver la entrega de VBUS y el retorno hacia la cámara. El contrato
de conexión deberá definir VBUS en el extremo anfitrión, secuencias
y prevención de retroalimentación. No se adopta como solución
desconectar el conductor de 5 V ni unir dos fuentes directamente.
La topología física de esa conexión permanece por concretar y no se
presenta como un cableado validado.

\subsubsection{Almacenamiento y mantenimiento}
\texttt{SD\_1} mantiene su asignación a la microSD, independiente
de los buses anteriores. Para mantenimiento se reserva
\texttt{USB\_1} como acceso de recuperación y \texttt{UART\_3}
como consola del procesador \cite{verdinds}. Sus niveles de referencia
deben respetarse; una referencia de tensión no implica alimentación
desde la herramienta externa. La asignación de los buses y de los conectores de señales se desarrolla
a continuación. Los GPIO de supervisión y las conexiones restantes
requieren comprobar multiplexación, estados al reiniciar y disipación.

\subsection{Asignación de buses y conectores de señales}
Los números de la tabla corresponden a contactos del conector X1 del
módulo Verdin, no a contactos del arnés externo ni a las bolas del SoC.
Se mantienen las funciones estándar documentadas por Toradex
\cite{verdinds}. Las redes externas se conectan mediante la adaptación
de niveles correspondiente; la tabla no autoriza conectar 3,3 V
directamente a señales de 1,8 V.

\begin{tabularx}{\textwidth}{L{3cm}L{2cm}X}
\toprule Función & Contacto X1 & Uso en la PCB \\
\midrule
I2C\_1\_SCL & 14 & Reloj de telemetría EPS. \\
I2C\_1\_SDA & 12 & Datos bidireccionales EPS. \\
UART\_1\_TXD & 131 & Transmisión hacia RX de actitud. \\
UART\_1\_RXD & 129 & Recepción desde TX de actitud. \\
SPI\_1\_MOSI & 200 & Datos hacia comunicaciones. \\
SPI\_1\_MISO & 198 & Datos desde comunicaciones. \\
SPI\_1\_CLK & 196 & Reloj generado por el computador. \\
SPI\_1\_CS & 202 & Selección del dispositivo de comunicaciones. \\
UART\_3\_TXD & 149 & Salida de consola. \\
UART\_3\_RXD & 147 & Entrada de consola. \\
\bottomrule
\end{tabularx}

Los contactos enumerados no se repiten. Esta comprobación de asignación
no sustituye la configuración de multiplexación y software: se deben
habilitar esas funciones y evitar sus alternativas en el mismo contacto.
No constituye todavía la asignación completa del módulo; faltan señales
de supervisión, alimentación, gestión de arranque y otros buses.

\subsubsection{Conexión al módulo de actitud}
Se utiliza un conector JST GH vertical de tres contactos
\texttt{BM03B-GHS-TBT}, con carcasa de cable \texttt{GHR-03V-S}
y contactos crimpados \texttt{SSHL-002T-P0.2}.
Se adopta conductor AWG 28 dentro del rango publicado de la familia.
La serie ofrece paso de 1,25 mm, polarización y bloqueo mecánico
\cite{jstgh}. Estas características permiten una conexión compacta y
reducen la posibilidad de inversión o desconexión accidental. Debe
comprobarse el espacio de inserción vertical y sujetarse el arnés para
no transmitir esfuerzos a las soldaduras.

\begin{tabularx}{\textwidth}{L{1.5cm}L{3cm}X}
\toprule Contacto & Red en J\_ADCS & Conexión y sentido \\
\midrule
1 & GND & Referencia común de señales. \\
2 & ADCS\_TX & Salida del computador a 3,3 V hacia RX externo. \\
3 & ADCS\_RX & Entrada al computador desde TX externo a 3,3 V. \\
\bottomrule
\end{tabularx}

La numeración es una decisión del diseño, no un estándar de pinout UART.
Debe conservarse en ambos extremos y verificarse con el dibujo del
fabricante, evitando invertir la vista del lado de acoplamiento.
El módulo de actitud recibe su energía externamente; este conector
transporta señales y referencia, sin alimentar actuadores.
El arnés tendrá inicialmente una longitud máxima de 150 mm, como
hipótesis geométrica para el diseño 1U que deberá verificarse en montaje.
Esta longitud no demuestra por sí sola integridad de señal.

\subsubsection{Conexión de telemetría al EPS}
Se separa J\_EPS en J\_EPS\_PWR y J\_EPS\_SIG.
La separación permite dimensionar contactos y cable de potencia a
partir de su corriente, sin imponer ese tamaño a las señales.
También permite comprobar la continuidad del arnés de telemetría por
separado. Esta división añade un conector y exige controlar la conexión
de señales cuando alguno de los extremos no está alimentado.

J\_EPS\_SIG utiliza \texttt{BM05B-GHS-TBT}, carcasa
\texttt{GHR-05V-S} y contactos \texttt{SSHL-002T-P0.2}, con cable
AWG 28. Emplear cinco contactos lo distingue físicamente del conector
de actitud de tres contactos. El uso de una misma familia simplifica
la preparación del arnés y de las huellas, sin hacer intercambiables
los conjuntos de diferente tamaño.

\begin{tabularx}{\textwidth}{L{1.5cm}L{3cm}X}
\toprule Contacto & Red & Función \\
\midrule
1 & GND & Referencia común de señales. \\
2 & EPS\_SCL & Reloj I2C, dominio externo de 3,3 V. \\
3 & EPS\_SDA & Datos I2C, dominio externo de 3,3 V. \\
4 & EPS\_PGOOD & Estado de potencia hacia supervisión local. \\
5 & EPS\_ALERT\_N & Aviso del EPS hacia el computador. \\
\bottomrule
\end{tabularx}

La asignación de PGOOD a la supervisión y de ALERT a un GPIO se cerrará
al diseñar el arranque y la adaptación. No se conectan arbitrariamente
a entradas del SoM. Se fija como hipótesis inicial un arnés de hasta
150 mm. La capacitancia del arnés, conectores, adaptador y pines se
incluirá al calcular los pull-ups; la longitud no sustituye ese cálculo.
El retorno de señales debe compartir referencia con potencia, evitando
que el conductor delgado del arnés se convierta en el retorno de carga
si se desconecta J\_EPS\_PWR. No se especifica conexión en caliente.

La ficha vigente de la familia GH publica un intervalo de
$-40$ a $+105\,^{\circ}\mathrm{C}$ y 1 A con conductor AWG 26
\cite{jstgh}. La corriente citada no se extrapola al AWG 28 ni se
utiliza para dimensionar la alimentación de la tarjeta. La selección
es para señales; tampoco acredita calificación espacial o resistencia
del conjunto frente a vibración. Antes de liberar la BOM deben
verificarse los sufijos de suministro, el plano y el rango térmico
de las referencias exactas adquiridas.

\subsubsection{Conexiones que requieren dimensionamiento adicional}
J\_EPS\_PWR se seleccionará después de sumar las corrientes por riel,
incluyendo transmisión de radio y simultaneidad. Los 3 A asignados al
SoM no representan la corriente total. La caída de tensión se evaluará
en el trayecto completo de ida y retorno, contactos, protecciones y
pistas; no se justificará capacidad mediante reparto ideal de corriente
entre contactos paralelos.
La cámara mantiene su enlace USB3 y requiere un conector apropiado
para esa interfaz. Su conjunto de alimentación externo debe cerrar
el tratamiento de VBUS. No se sustituye por un conector GH.
El SPI permanece como pistas internas entre las dos secciones de la
misma PCB. No se añade un conector entre computador y comunicaciones.

\subsection{Adaptación de niveles de la UART de actitud}
La UART se implementa mediante el TXU0202DCUR, con un canal desde el
computador hacia actitud y otro en sentido opuesto. La dirección fija
corresponde a la función TX/RX y evita controles de dirección adicionales.
La alternativa SN74AXC2T245 admite direcciones configurables y umbrales
especificados por intervalos de alimentación; se conserva como alternativa
si la revisión de márgenes exige modificar el traductor
\cite{txu0202,axc2t245}. La selección actual permite la captura inicial,
pero no se considera una validación completa del bloque.

\subsubsection{GPIO de habilitación y estados de arranque}
Se asigna GPIO\_1, contacto X1.206 del Verdin, a
\texttt{UART\_ADCS\_EN}. Corresponde a GPIO1\_IO00 del SoC y a la
función GPIO1\_IO0, ALT0. Es un GPIO dedicado y no comparte los
contactos asignados a UART\_1, I2C\_1 o SPI\_1. La tabla 17 del
Verdin indica polarización hacia bajo en el estado de reinicio y la
tabla 21 identifica su función GPIO \cite{verdinds}.

La red se conecta a OE, pin 6 del TXU0202, y lleva R\_OE de
10 k$\Omega$ a GND. Un nivel bajo mantiene las salidas deshabilitadas.
El estado de reinicio documentado no equivale a una garantía de que
el cargador de arranque y Linux conservarán ese estado: ambos deberán
mantener el GPIO bajo hasta configurar TX en reposo alto y verificar
las alimentaciones. La implementación de software queda fuera de esta
práctica, pero esta condición se registra como requisito de integración.
La resistencia no puede contrarrestar una salida GPIO activamente alta.

La secuencia prevista es: GPIO bajo; estabilización de fuentes;
configuración de UART y desactivación de polarizaciones internas que
interfieran con las externas; TX en reposo alto; habilitación del enlace.
Antes del apagado controlado se deshabilita OE. En reinicio se vuelve
al estado de deshabilitación. Los datos recibidos durante una transición
se descartan; una transición eléctrica no debe interpretarse como orden
válida de actuación. No se incluye suspensión profunda como modo
operativo validado de este bloque.

\subsubsection{Origen de los 1,8 V y presupuesto de alimentación}
VCCA, pin 3, se conecta a \texttt{PWR\_1V8\_MOCI}, contacto X1.214.
La salida documentada tiene capacidad máxima de 250 mA. La guía Verdin
recomienda utilizar este dominio para circuitos conectados a sus señales,
por su relación con las fuentes de E/S y la prevención de alimentación
involuntaria \cite{verdinds,verdinguide}. No se conecta a una fuente
externa de 1,8 V ni se confunde con una entrada de alimentación del SoM.

VCCB, pin 7, se conecta al riel local de 3,3 V de la tarjeta.
El GND del traductor, el del SoM y J\_ADCS.1 comparten referencia.
Se coloca un condensador de 100 nF entre cada alimentación y GND,
próximo al contacto correspondiente \cite{txu0202}.

Se reserva inicialmente un presupuesto de diseño de 1 mA del riel
PWR\_1V8\_MOCI para el conjunto UART. Equivale al 0,4\,\% de los
250 mA publicados; no es un consumo máximo demostrado ni autoriza
ignorar las demás cargas de esa salida. La suma de cargas y los
transitorios se verificarán en el presupuesto general. La capacidad
de la fuente y la demanda real se mantienen como magnitudes distintas.

Para el análisis se conservan envolventes de 1,65--1,95 V y 3,0--3,6 V.
Son límites que deberá satisfacer el diseño, no tolerancias atribuidas
sin comprobación a la salida de Toradex. Si los rieles o las caídas
reales exceden estas envolventes, se recalculan los márgenes.

\subsubsection{Conexiones completas del traductor}
Los números corresponden exclusivamente al encapsulado DCU.
Los designadores funcionales se convertirán en designadores únicos
al anotar el proyecto Altium.
\begin{tabularx}{\textwidth}{L{1cm}L{1.8cm}X}
\toprule Pin & Nombre & Conexión \\
\midrule
1 & B2 & J\_ADCS.3, ADCS\_RX; pull-up de 10 k$\Omega$ a 3,3 V. \\
2 & GND & Referencia común y retornos de desacoplo. \\
3 & VCCA & X1.214, PWR\_1V8\_MOCI; 100 nF a GND. \\
4 & A2Y & X1.129, UART\_1\_RXD; pull-up de 47 k$\Omega$ a VCCA. \\
5 & A1 & X1.131, UART\_1\_TXD; pull-up de 10 k$\Omega$ a VCCA. \\
6 & OE & X1.206, GPIO\_1; pull-down de 10 k$\Omega$ a GND. \\
7 & VCCB & Riel local de 3,3 V; 100 nF a GND. \\
8 & B1Y & J\_ADCS.2, ADCS\_TX. \\
\bottomrule
\end{tabularx}

Las resistencias definen estados en alta impedancia; no son
terminaciones de línea ni adaptadores de tensión. Durante el reinicio,
las polarizaciones internas de los pines UART pueden competir con las
externas. No se supone reposo alto garantizado hasta configurar los
pines; OE permanece bajo durante esa etapa.

\subsubsection{Contrato eléctrico del módulo externo}
El módulo de actitud recibe su alimentación externamente. Se exige
que su salida TX entregue al menos 2,9 V en alto y como máximo 0,2 V
en bajo, con carga de hasta 0,5 mA. Se exige a su entrada RX reconocer
alto a partir de 2,52 V y bajo hasta 0,90 V. Son especificaciones
impuestas al módulo externo idealizado, no resultados de un módulo
comercial ensayado.

La entrada RX externa llevará polarización de 47 k$\Omega$, 1\,\%,
a su propio riel, limitado a 3,6 V. Su fuga será como máximo
10 $\mu$A por pin, incluyendo la condición de módulo apagado; sus
pines tolerarán señales hasta 3,6 V sin alimentar su fuente interna.
Se conserva el arnés de hasta 150 mm y la referencia GND común.
El límite de 10 $\mu$A no incluye una resistencia que conecte
intencionadamente la línea a un riel apagado de baja impedancia;
el módulo externo debe resolver ese caso para cumplir el contrato.

La protección de apagado del traductor local no detecta por sí sola
que el ADCS remoto se ha apagado mientras VCCA y VCCB siguen activas.
Por ello el contrato de alta impedancia es obligatorio. Si no se cumple,
se modifica el aislamiento externo o la interfaz. No se especifica
conexión en caliente.

\subsubsection{Cálculos y márgenes: resultados y límites}
Para las resistencias seleccionadas se consideran tolerancia inicial
del 1\,\% y coeficiente térmico de 100 ppm/$^{\circ}$C.
En $-40$ a $+85\,^{\circ}$C, la desviación máxima respecto a
$25\,^{\circ}$C es 65 $^{\circ}$C. Como envolvente conservadora:
\begin{align}
R_{\min}&=R_n(1-0,01)(1-0,0065)=0,983565R_n,\\
R_{\max}&=R_n(1+0,01)(1+0,0065)=1,016565R_n.
\end{align}
Para 47 k$\Omega$: $R_{\min}=46,228$ k$\Omega$ y
$R_{\max}=47,779$ k$\Omega$. Para 10 k$\Omega$:
$R_{\min}=9,836$ k$\Omega$ y $R_{\max}=10,166$ k$\Omega$.
Estas cotas no incluyen envejecimiento; una deriva adicional deberá
incorporarse si se convierte en requisito de diseño.

El cambio de 10 k$\Omega$ a 47 k$\Omega$ en A2Y reduce la corriente
que debe absorber el traductor cuando RX está en bajo. Usando la fuga
de 4 $\mu$A de la tabla GPIO de NXP como cota conservadora y sin
polarización interna adicional habilitada \cite{imx8mpiec}:
\begin{equation}
I_{A2Y,L}\leq\frac{1,95}{46227,555}+4\ \mu\mathrm{A}
=46,18\ \mu\mathrm{A}<100\ \mu\mathrm{A}.
\end{equation}
Así se puede utilizar la especificación del TXU0202 para carga de
100 $\mu$A: $V_{OH}\geq V_{CCO}-0,1$ V y $V_{OL}\leq0,1$ V.
No se utiliza esta especificación con una carga mayor.
Con $V_{IH}=0,7V_{DD}$ y $V_{IL}=0,3V_{DD}$ del SoC,
y evaluando extremos independientes de la envolvente:
\begin{align}
NM_H&\geq(1,65-0,1)-0,7(1,95)=0,185\ \mathrm{V},\\
NM_L&\geq0,3(1,65)-0,1=0,395\ \mathrm{V}.
\end{align}
Estos resultados corresponden al recorrido traductor hacia RX del
Verdin y a las condiciones de carga indicadas; no validan los demás
recorridos automáticamente.

Para B1Y se limita la carga estática externa mediante el contrato:
\begin{equation}
I_{B1Y,L}\leq\frac{3,6}{46227,555}+10\ \mu\mathrm{A}
+0,05\ \mu\mathrm{A}=87,93\ \mu\mathrm{A}<100\ \mu\mathrm{A}.
\end{equation}
El último término corresponde a la envolvente de fuga del TVS a su
tensión de trabajo. Resultan $NM_H\geq2,9-2,52=0,38$ V y
$NM_L\geq0,90-0,1=0,80$ V hacia el receptor externo. El margen
se consume con ruido, diferencias de tierra y caída en el arnés;
no se interpreta como una tensión de ruido permitida sin presupuesto.

La dirección SoM hacia A1/OE necesita especial atención. NXP especifica
$V_{OH}\geq0,8V_{DD}$ y $V_{OL}\leq0,2V_{DD}$ en las condiciones
de carga de su tabla GPIO. TI tabula para entrada de 1,65 V los
extremos Schmitt $V_{T+\max}=1,13$ V y $V_{T-\min}=0,35$ V.
En ese punto documentado, con ambos dominios a 1,65 V:
\begin{align}
NM_H&\geq0,8(1,65)-1,13=0,190\ \mathrm{V},\\
NM_L&\geq0,35-0,2(1,65)=0,020\ \mathrm{V}.
\end{align}
El margen bajo es reducido. No se extrapola esa fila de TI a toda
la banda 1,65--1,95 V ni se interpolan umbrales como si fueran límites
garantizados. Por tanto, no se declara cerrado el margen de A1/OE
en toda la envolvente. Para liberarlo se requiere información de
fabricante que cubra el intervalo o sustituir el traductor por una
solución con límites adecuados y explícitos. Una medición típica puede
apoyar la revisión, pero no demuestra un máximo garantizado.

Análogamente, los niveles externos propuestos ofrecen márgenes de
0,98 V y 0,69 V hacia B2 en el punto TI de 3,0 V
($V_{T+\max}=1,92$ V; $V_{T-\min}=0,89$ V). No se atribuyen
estos márgenes a toda la banda 3,0--3,6 V sin comprobarla.

En OE deshabilitado, la cota de fuga de 4 $\mu$A del SoC más
2 $\mu$A del traductor produce:
\begin{equation}
V_{OE}\leq6\ \mu\mathrm{A}\cdot10165,65\ \Omega
=0,0610\ \mathrm{V}.
\end{equation}
La polarización interna hacia bajo del GPIO ayuda en reinicio. Esta
cota supone GPIO en alta impedancia; no describe una salida activa.

Para los pull-ups de 10 k$\Omega$, las nuevas cotas térmicas dan
0,1983 mA y 0,3660 mA a 1,95 V y 3,6 V, respectivamente. La mayor
disipación por resistencia de este grupo es 1,318 mW. Para el pull-up
local de 47 k$\Omega$ a 1,95 V resulta 0,0823 mW.
El presupuesto estático conservador del lado A, sumando tres ramas
de polarización y reservando 6 $\mu$A para el traductor, es inferior
a 0,45 mA. La reserva de 1 mA no queda demostrada para todo régimen
dinámico por este cálculo estático.

\subsubsection{Referencias de soporte y protección}
Se proponen las siguientes referencias para captura y preparación de
librerías; esta lista no sustituye la BOM exportada desde Altium.
\begin{longtable}{L{2.5cm}L{1cm}L{5cm}L{4cm}}
\toprule Función & Cant. & Referencia & Especificación \\
\midrule\endhead
Traductor & 1 & TXU0202DCUR & DCU, ocho contactos. \\
Pull-up A1 y B2; pull-down OE & 3 & RC0603FR-0710KL & 10 k$\Omega$, 1\,\%, 0603. \\
Pull-up A2Y & 1 & RC0603FR-0747KL & 47 k$\Omega$, 1\,\%, 0603. \\
Desacoplo A y B & 2 & C0603C104K3RACTU & 100 nF, 25 V, 10\,\%, X7R, 0603. \\
Protección primaria TX y RX & 2 & ESD321DPYR & DPY, TVS a GND. \\
\bottomrule
\end{longtable}

Las resistencias RC tienen potencia nominal de 0,1 W a 70 $^{\circ}$C;
su curva de reducción térmica debe asociarse a la librería/BOM al
liberarla \cite{yageo10k,yageo47k}. Las disipaciones calculadas son
pequeñas respecto a esa potencia, pero no se afirma que conserve
0,1 W a cualquier temperatura. Se elige 0603 para facilitar montaje
e inspección y porque no existe aquí una exigencia de miniaturización
que justifique un formato menor.

Los condensadores de 25 V sustituyen la especificación inicial de
10 V sin cambiar su función ni su valor nominal. La tensión de trabajo
de 3,6 V representa el 14,4\,\% de su tensión nominal. Esto no
garantiza por sí solo la capacitancia efectiva: tolerancia, temperatura,
sesgo DC y envejecimiento deben evaluarse. Considerando únicamente
10\,\% de tolerancia y 15\,\% térmico, la cota sería 76,5 nF;
no se presenta como mínimo total garantizado \cite{kemet100n}.

Se propone un ESD321DPYR por línea, inmediatamente junto al conector:
pin 1 a la señal y pin 2 a GND. El pinout corresponde a DPY; no se
debe reutilizar el de DYA, que tiene numeración diferente.
El retorno será corto y directo al plano de referencia, evitando
desviar la descarga a través del SoM. La fuga de trabajo y la baja
capacitancia son compatibles con el presupuesto estático anterior
\cite{esd321}. No se conectan estos TVS al riel de 1,8 V.

Esta selección es protección primaria propuesta, no inmunidad ESD
validada. El fabricante publica una tensión de limitación típica de 6,8 V a
16 A TLP, mientras que el máximo absoluto de entrada del TXU0202 es
6,5 V. Por ello no se afirma que el TVS mantenga siempre el pin dentro
de ese límite. Tampoco se interpreta la diferencia como una prueba
automática de fallo bajo ESD: son condiciones de ensayo distintas.
Antes de liberar la protección deben coordinarse las curvas de pulso,
el circuito secundario si resulta necesario y las inductancias del
montaje. No se asigna una clasificación IEC al conjunto a partir de
la clasificación del diodo.

\subsubsection{Paso al esquemático y criterio de cierre}
Puede iniciarse la captura del bloque en la hoja del computador,
incluyendo conector, traductor, polarizaciones, desacoplo y puntos de
prueba de VCCA, VCCB, OE, TX, RX y GND. El conector de actitud no
transporta potencia de actuadores. El bloque permanece identificado
como en revisión hasta resolver los márgenes de entrada y la
coordinación ESD anteriores.

Se revisarán además numeración símbolo--huella, estados de arranque,
ausencia de conflictos de multiplexación y conectividad de tierras.
El ERC y la revisión manual se ejecutarán sobre el circuito capturado.
Comenzar la captura no implica que el informe completo esté terminado:
las demás interfaces, alimentación de tarjeta, memoria, restricciones
de fabricación, BOM y costos conservan sus etapas de desarrollo.

\subsection{Necesidades de procesamiento y memoria}
La cámara de referencia es la XIMEA
\texttt{MQ022HG-IM-SM5X5-7NIR2}. Se adopta para el diseño académico un
módulo Verdin iMX8M Plus Quad 4GB IT como computador principal, con RAM
y memoria de arranque incorporadas. El almacenamiento persistente de
imágenes, telemetría y resultados se implementa por separado. Los
archivos se conservan en ese almacenamiento y se cargan en RAM cuando
se procesan. La selección se desarrolla en la sección de componentes.

El alcance de cómputo considerado comprende adquisición RAW16, gestión
de archivos y metadatos, preparación de datos para comunicaciones y
supervisión de estados. No se presupone un algoritmo de inteligencia
artificial ni se acredita procesamiento espectral avanzado o desempeño
de control en tiempo real. Sus recursos y plazos deberán comprobarse
si esas funciones se incorporan. Los presupuestos heredados se revisan
con las características de los dispositivos elegidos.


La evaluación del procesador deberá comprobar:
\begin{itemize}
\item Periféricos compatibles y simultáneos para cámara, memoria, radio, potencia y actitud.
\item Ausencia de conflictos de multiplexación de pines en la configuración completa.
\item Capacidad de transferir datos y atender interrupciones sin perder información.
\item Memoria de programa, RAM y buffers suficientes; DMA cuando resulte necesario para las tasas acordadas.
\item Recursos de arranque, reinicio, supervisión y programación, con estados definidos de las salidas.
\end{itemize}

Para un conjunto de $N_b$ imágenes de banda, cada una con $N_x\times N_y$ muestras y una profundidad de $b$ bits por muestra, el volumen ideal sin compresión es
\begin{equation}
D_{\mathrm{img}}=\frac{N_xN_yN_b b}{8}\quad[\mathrm{bytes}],
\end{equation}
Esta expresión supone muestras empaquetadas sin relleno. Si cada muestra
ocupa un número entero de bytes, el factor $b/8$ se sustituye por
$\lceil b/8\rceil$. El volumen real incluirá el relleno por línea y los
metadatos del formato seleccionado.

En una cámara con filtros espectrales en mosaico, un cuadro bruto contiene
una muestra por posición del sensor, no una muestra de cada banda en cada
posición. Si el sensor tiene $N_{x,s}\times N_{y,s}$ posiciones, el volumen
ideal del cuadro bruto es
\begin{equation}
D_{\mathrm{RAW}}=\frac{N_{x,s}N_{y,s}b}{8}\quad[\mathrm{bytes}],
\end{equation}
con las mismas correcciones de empaquetado. El volumen reconstruido se
calculará con las dimensiones reales de las imágenes de banda resultantes.
Se definirá qué representación se almacena y transmite; no se sumarán
ambas salvo que el sistema deba conservarlas simultáneamente.

Si la arquitectura permite almacenar por bloques, la RAM no necesita contener una imagen completa. Como criterio de dimensionamiento se utilizará
\begin{equation}
M_{\mathrm{RAM}}\geq M_{\mathrm{base}}+M_{\mathrm{adq}}+M_{\mathrm{alm}}+M_{\mathrm{radio}}+M_{\mathrm{reserva}}.
\end{equation}
Los términos representan uso base del software previsto, buffers de adquisición, almacenamiento y radio, y reserva justificada. Su simultaneidad debe comprobarse; no se asignan valores sin conocer los buses y los modos de operación. Dos buffers de adquisición, por ejemplo, solo se justificarían si se necesita recibir un bloque mientras se procesa o almacena el anterior.

Para el almacenamiento de misión se utilizará el siguiente criterio conservador de retención sin descarga:
\begin{equation}
C_{\mathrm{util}}\geq C_0+\dot{D}_{\mathrm{total}}T_{\mathrm{ret}}+C_{\mathrm{reserva}},
\end{equation}
donde $C_0$ es el volumen ya almacenado al comenzar el intervalo,
$\dot{D}_{\mathrm{total}}$ es la generación prevista en bytes por día,
incluyendo imágenes, telemetría y eventos, y $T_{\mathrm{ret}}$ se expresa
en días. Así se evita contabilizar los eventos dos veces. La reserva es
espacio libre adicional previsto para operación y contingencias; la
capacidad útil excluye el espacio no disponible para datos de misión,
como el ocupado por el formato o el software. ET-04 se conserva hasta
justificar y documentar cualquier modificación.

\subsection{Compatibilidad entre datos y comunicaciones}
\label{sec:datos}
El presupuesto de datos vincula la capacidad de adquisición, el almacenamiento y el enlace de comunicaciones. Estas funciones deben dimensionarse conjuntamente para evitar acumulación sostenida de información sin posibilidad de descarga.

\subsubsection{Comprobación del presupuesto heredado}
El cuadro 19 de la práctica 3 \cite{lab3} estima 10 adquisiciones diarias de 5 MB cada una. Adoptando MB decimales y 600 s diarios de contacto como en ese presupuesto:
\begin{align}
D_{\mathrm{imagenes,dia}}&=10(5\times10^6)=50\times10^6\ \mathrm{bytes},\\
R_{\mathrm{imagenes,util}}&=\frac{8(50\times10^6)}{600}
\approx666\,667\ \mathrm{bit/s}\approx667\ \mathrm{kbit/s}.
\end{align}
Al añadir la estimación heredada de 1,2 MB diarios de telemetría, la tasa útil total es aproximadamente 683 kbit/s. La tasa física deberá ser mayor si existen cabeceras, corrección de errores, retransmisiones o periodos no aprovechables. Si $\eta$ es la fracción útil del enlace,
\begin{equation}
R_{\mathrm{fisica}}\geq\frac{8D_{\mathrm{total,dia}}}{\eta\,T_{\mathrm{contacto}}},\qquad 0<\eta\leq1.
\end{equation}
En esta expresión, $D_{\mathrm{total,dia}}$ es el volumen generado durante un día y $T_{\mathrm{contacto}}$ es el tiempo de contacto asignado a su descarga. La fracción $\eta$ incluye la sobrecarga y el aprovechamiento efectivo dentro de ese intervalo; esos factores no deben descontarse nuevamente del tiempo. Su valor se determinará con el protocolo y el escenario de operación. Los 600 s son un supuesto heredado, no un resultado de análisis orbital realizado en esta práctica.

Si se eligiera operar a 9600 bit/s, el máximo ideal durante 600 s sería
\begin{equation}
D_{9600}=\frac{9600(600)}{8}=720\,000\ \mathrm{bytes}=0,72\ \mathrm{MB}.
\end{equation}
Por tanto, una operación a esa tasa no descarga el volumen estimado, incluso sin sobrecarga. ET-06 especifica un mínimo y permite escoger una tasa mayor; la contradicción aparece al tratar 9600 bit/s como solución suficiente. Incrementar la memoria solo retrasa la saturación si la generación supera permanentemente la descarga.

\subsubsection{Criterio adoptado para el nuevo presupuesto}
El presupuesto de datos se recalculará a partir de las características
de la cámara seleccionada. Se mantiene como escenario preliminar una
frecuencia de diez adquisiciones diarias y un tiempo total de contacto
de 600 segundos por día. Ambos valores se consideran supuestos de
dimensionamiento sujetos a revisión; el tiempo de contacto deberá
comprobarse mediante el análisis de las oportunidades de comunicación.

El volumen de cada adquisición se determinará con la resolución, el
formato de salida, el empaquetado de píxeles, los metadatos y las bandas
efectivamente adquiridas y conservadas. En sensores con filtros
espectrales en mosaico se distinguirá el tamaño del cuadro bruto del
volumen de las imágenes espectrales reconstruidas, evitando contabilizar
varias veces los mismos datos. El cálculo inicial utilizará datos sin
compresión, de modo que la viabilidad del enlace no dependa de una
reducción de tamaño todavía no demostrada.

La telemetría se dimensionará por separado a partir de las variables
registradas, los bytes por registro y la frecuencia de generación.
En consecuencia, los 50 MB diarios de imágenes y los 1,2 MB diarios
de telemetría de la práctica 3 se conservan únicamente como antecedentes
del análisis anterior y serán sustituidos por los resultados de este
dimensionamiento.

El volumen diario resultante y el tiempo de contacto establecerán la
tasa útil requerida del enlace. La selección de comunicaciones deberá
considerar además la sobrecarga del protocolo y el margen necesario
para las condiciones de operación previstas. Si la capacidad de descarga
resulta insuficiente, se revisarán conjuntamente la frecuencia de
adquisición, el volumen conservado por captura y la capacidad del enlace.
El computador seleccionado y la elección del radio deberán ser compatibles
con este presupuesto y con las tasas de transferencia internas.

\subsection{Alimentación y presupuesto de consumo}
\subsubsection{Antecedentes y criterio de dimensionamiento}
El cuadro 18 de la práctica 3 \cite{lab3} asigna a OBC y memoria 150 mW nominales y 300 mW máximos. Son estimaciones, no resultados de selección. Para observar su equivalencia a corriente, suponiendo que toda esa potencia se entrega en un único riel ideal de 3,3 V:
\begin{equation}
I_{\mathrm{nom}}=\frac{0,150}{3,3}\approx45,5\ \mathrm{mA},\qquad
I_{\mathrm{max}}=\frac{0,300}{3,3}\approx90,9\ \mathrm{mA}.
\end{equation}
Estas corrientes son equivalencias de las estimaciones heredadas y no dimensionan todavía el conector ni constituyen un límite comprobado. Dado que la selección admite un procesador de aplicaciones o un módulo de cómputo, el presupuesto de 300 mW no se adopta como límite validado de la solución. El supuesto de un solo riel excluye pérdidas y debe sustituirse por el reparto real de consumos. La memoria puede presentar picos de escritura que no estén representados en el promedio heredado.

El subsistema externo de potencia se modela como una fuente que satisface
las condiciones de entrada especificadas por este diseño. La estimación
global heredada de 5 W no se impone como límite del diseño actualizado.
Para la rama dedicada al computador se establece una capacidad de
suministro de 3 A a 5 V, manteniendo una tolerancia de
$\pm5\,\%$ en los pines del módulo. Esta capacidad no representa el
consumo permanente ni incluye cámara, almacenamiento externo,
comunicaciones o auxiliares. El requerimiento total de la PCB se
obtendrá al completar esos bloques y sus pérdidas.

Se especifica una tensión de entrada de $5\ \mathrm{V}\pm2\,\%$
entre los nodos \texttt{+5V\_EPS\_IN} y GND de la PCB, después
del conector de alimentación. El EPS debe satisfacer este contrato
considerando las caídas anteriores a dicha frontera. La rama del SoM
incorpora una protección independiente entre esta entrada y
\texttt{+5V\_SOM}; su componente y ajustes requieren verificación
eléctrica antes de cerrar la selección.
El margen inferior de distribución a 3 A es
\begin{equation}
\Delta V_{\mathrm{distrib}}=4{,}90-4{,}75=0{,}15\ \mathrm{V},
\qquad
R_{\mathrm{trayecto,max}}=\frac{0{,}15}{3}
=50\ \mathrm{m}\Omega.
\end{equation}
Esta resistencia equivalente incluye protección, distribución y
retorno desde la frontera definida hasta los pines del módulo.
El diseño debe quedar por debajo de esa cota para reservar margen
frente a transitorios; consumirla por completo solo satisface el
límite estático. La tensión especificada no sustituye la comprobación
de arranque, sobreimpulso y respuesta a cambios de carga.

Para cada riel se calculará la corriente máxima aplicable sumando dispositivos activos en cada modo y sus circuitos auxiliares. Si existe conversión local, su corriente de entrada se estimará mediante
\begin{equation}
I_{\mathrm{entrada}}=\frac{\sum_k V_k I_k}{\eta_{\mathrm{reg}}V_{\mathrm{entrada}}},
\end{equation}
Esta expresión supone una eficiencia equivalente para el conjunto de salidas. Si se emplean convertidores independientes, se calculará la potencia de entrada de cada uno con su eficiencia y carga, y se sumarán sus contribuciones junto con las cargas directas. El consumo propio solo se añadirá por separado cuando no esté incluido en la eficiencia utilizada, evitando contabilizarlo dos veces.


\subsubsection{Desacoplo y evaluación de la protección del computador}
La alimentación del SoM utiliza los contactos 251, 253, 255, 257 y
259 de X1, unidos a \texttt{+5V\_SOM}; los contactos GND se conectan
al retorno común. Se incorpora un banco en paralelo de cinco
condensadores de 4,7 $\mu$F, 10 V, X5R, tolerancia del 10\,\% y
encapsulado 0603, más uno de 22 $\mu$F, 10 V, X5R, tolerancia del
20\,\% y encapsulado 0805. Este conjunto toma como base el circuito
de referencia de Toradex, que indica su colocación próxima a los
contactos de alimentación del zócalo \cite{verdinguide}.
\begin{align}
C_{\mathrm{nom}}&=5(4{,}7)+22=45{,}5\ \mu\mathrm{F},\\
C_{\min,\mathrm{tol}}&=5(4{,}7)(0{,}90)+22(0{,}80)
=38{,}75\ \mu\mathrm{F},\\
C_{\max,\mathrm{tol}}&=5(4{,}7)(1{,}10)+22(1{,}20)
=52{,}25\ \mu\mathrm{F}.
\end{align}
Los extremos anteriores incluyen únicamente tolerancia inicial;
no son límites garantizados de capacidad efectiva en servicio.
La selección comercial debe comprobar polarización continua,
temperatura y envejecimiento. A 5,25 V se utiliza el 52,5\,\% de
la tensión nominal de 10 V; esta relación no cuantifica la pérdida
de capacitancia. La verificación de arranque debe incluir también
la capacidad y la carga internas del módulo.

Se evalúan dispositivos eFuse para la protección independiente de
la rama. A 3 A, el TPS25947 especifica una resistencia máxima de
45 m$\Omega$ para la envolvente de entrada y temperatura indicada
en su tabla eléctrica \cite{tps25947}. Por tanto,
\begin{equation}
\Delta V_{47}=3(0{,}045)=0{,}135\ \mathrm{V},\qquad
P_{47}=3^2(0{,}045)=0{,}405\ \mathrm{W}.
\end{equation}
Solo quedarían 15 mV del presupuesto de 150 mV para distribución,
retorno y transitorios. No se considera cerrada esta selección.
El TPS25948 ofrece una resistencia máxima de 20 m$\Omega$ a 3 A
en las condiciones de su tabla eléctrica \cite{tps25948}, con
caída de 60 mV y pérdidas de conducción de 0,180 W. Esta alternativa
deja 90 mV de margen, pero requiere revisar su encapsulado PowerWCSP,
la protección ante entrada negativa, la variante y los ajustes.
El bloqueo de corriente inversa no demuestra protección ante
polaridad invertida. Estos cálculos comparan conducción estática;
no validan el comportamiento de arranque, sobrecorriente o térmico
del circuito completo.

\subsubsection{Condiciones de arranque, reinicio y supervisión}
El Verdin iMX8M Plus incorpora el circuito que inicia el encendido
al aplicar la alimentación principal \cite{verdinds}. Para servicio
se incorporan dos pulsadores normalmente abiertos a GND: uno en
X1.246, red \texttt{OBC\_RECOVERY\_N}, y otro en X1.260, red
\texttt{OBC\_RESET\_IN\_N}. El primero permite solicitar recuperación
durante el encendido y el segundo reiniciar el módulo. Las entradas
incluyen pull-ups internos de 10 k$\Omega$ y 100 k$\Omega$,
respectivamente, por lo que no se añaden pull-ups externos a 3,3 V
o 5 V \cite{verdinds}. La captura utiliza el símbolo de referencia
\texttt{SW-WMF-TD-03XB}; la correspondencia física de sus contactos
y la referencia comercial deben verificarse antes de liberar la BOM.
El pin X1.249, \texttt{VCC\_BACKUP}, se mantiene separado de la
alimentación principal mientras se define el respaldo del RTC.

Se establecen condiciones de arranque y reinicio que eviten operaciones
involuntarias mientras la alimentación se estabiliza y el computador
inicializa sus periféricos. El computador deberá permanecer en reinicio
hasta que su alimentación y las condiciones de arranque cumplan los
requisitos del dispositivo seleccionado. Asimismo, se definirá su
respuesta ante una caída de tensión fuera del intervalo de operación.

Durante el arranque, la transmisión de radio permanecerá deshabilitada
hasta recibir una autorización explícita del sistema, sin que ello
implique necesariamente desactivar la recepción. La orden de captura
de la cámara permanecerá inactiva hasta completar su alimentación y
configuración. Para el almacenamiento se establecerán los estados
inactivos de las señales de selección o escritura que correspondan a
la interfaz elegida. Las salidas hacia los controladores de actitud
adoptarán estados iniciales acordados con ese subsistema para evitar
activaciones involuntarias.

La implementación de estas condiciones se determinará con las hojas de
datos de los componentes seleccionados. Se comprobarán los estados de
los pines durante el encendido y el reinicio, los umbrales y tiempos de
supervisión y las secuencias de alimentación. Cuando resulte necesario,
se incorporarán resistencias de polarización, señales de habilitación
o un supervisor externo. Su selección se justificará considerando las
funciones ya integradas en cada dispositivo y el comportamiento cuando
los bloques conectados no estén alimentados simultáneamente.

La prevención de operaciones involuntarias durante el arranque se
distingue de la protección de datos frente a una interrupción de energía
durante una escritura. Esta última condición se evaluará con el
almacenamiento seleccionado, su comportamiento ante pérdida de
alimentación y los recursos disponibles para detectar y gestionar el
evento. Por tanto, no se considerará resuelta únicamente mediante la
polarización de señales o el reinicio del computador.

\subsection{Espacio y condiciones de integración}
El máximo preliminar de $90\times90$ mm corresponde a toda la PCB. No se establece un reparto porcentual entre central y radio sin conocer encapsulados y restricciones. El cierre mecánico requiere dimensiones útiles, agujeros, zonas de exclusión, altura admisible, posición de conectores y espacio para acceso de depuración.

La integración de digital y RF obliga a acordar zonas funcionales y recorridos de corriente. No se prescribe partir el plano de tierra sin análisis: la continuidad de retornos y las recomendaciones de los componentes se revisarán en los puntos de restricciones. El rango industrial de temperatura, por sí solo, no acredita aptitud espacial; las simplificaciones académicas se conservarán explícitas.

\subsection{Registro inicial de cambios}
\begin{tabularx}{\textwidth}{L{1.1cm}XX}
\toprule ID & Cambio o criterio adoptado & Justificación y efecto \\
\midrule
C-01 & Integrar central y comunicaciones en una PCB. & Actualización de la partición física asignada. IF-10/11 se vuelven internas; se comparte área y presupuesto eléctrico. \\
C-02 & Corregir método de verificación de ET-07. & Usar condiciones de operación recomendadas en lugar de máximos absolutos, en coherencia con la guía 4 \cite{guia4}. \\
C-03 & Complementar ET-08 con recursos y tasas. & La frecuencia no demuestra compatibilidad de cámara, memoria y periféricos. \\
C-04 & Revisar presupuesto de datos frente al enlace. & 9600 bit/s no satisface el volumen diario supuesto. Puede afectar adquisición, almacenamiento y comunicaciones. \\
C-05 & Precisar RNF-03 según la interfaz elegida. & Control de impedancia y pares diferenciales solo se aplican cuando la interfaz y sus condiciones lo requieren. \\
C-06 & Adoptar cámara XIMEA y computador modular Verdin. & Se separa el almacenamiento de misión de la RAM y memoria de arranque. Se revisan potencia e interfaces; la compatibilidad del SDK en el módulo permanece como hipótesis de integración. \\
C-07 & Adoptar requisitos de arranque y supervisión. & Definir estados iniciales y respuesta ante pérdida de alimentación; los circuitos concretos se seleccionarán posteriormente. \\
C-08 & Seleccionar almacenamiento independiente de misión. & microSD industrial de 16 GB y escenario de 90 días. Se distinguen capacidad nominal, umbral útil, cuotas y reserva. \\
C-09 & Definir contratos de módulos externos. & EPS por I2C y controlador ADCS por UART, con dominio externo de 3,3 V. La energía de cámara se entrega externamente; la topología VBUS requiere cierre eléctrico. \\
C-10 & Asignar buses y conectores de señales. & Contactos X1 para I2C, UART y SPI; conectores GH de tres y cinco contactos para actitud y telemetría EPS. La entrada de potencia se dimensiona por separado. \\
C-11 & Adaptar la UART de actitud. & TXU0202DCUR, polarización y habilitación controlada. Se exige tolerancia de señales con el módulo externo apagado y se mantienen verificaciones de umbrales y protección. \\
\bottomrule
\end{tabularx}

\section{Selección de componentes}
\subsection{Cámara de referencia para el dimensionamiento de la tarjeta}

Se selecciona la cámara XIMEA \texttt{MQ022HG-IM-SM5X5-7NIR2}
como referencia eléctrica de la carga útil académica. La cámara es un
dispositivo externo a la PCB central y de comunicaciones; su interfaz
y sus necesidades de alimentación condicionan el diseño de la tarjeta,
pero su costo no se incluye en la BOM de componentes montados en ella.

El modelo seleccionado dispone de un sensor de $2048\times1088$ píxeles,
conversión de 10 bits y conexión USB3 \cite{ximea7nir2}. La tarjeta
requerirá un computador con función anfitrión y software compatible con
la cámara. La disponibilidad de soporte Linux ARM orienta la selección,
pero deberá comprobarse en la plataforma concreta \cite{ximealinux}.

Se adopta para el dimensionamiento preliminar una adquisición de un
cuadro bruto completo en RAW16. Este formato conserva las muestras en
contenedores de 16 bits. El manual recomienda RAW8 o RAW16 para obtener
los datos originales \cite{ximeamanual}. Sin metadatos ni relleno
adicional, el volumen por cuadro es
\begin{equation}
D_{\mathrm{cuadro}}=2048\times1088\times2
=4\,456\,448\ \mathrm{bytes}.
\end{equation}
Para diez adquisiciones diarias se obtienen $44{,}56448$ MB por día,
utilizando MB decimales. La descarga de esas imágenes en 600 s requiere
aproximadamente $594{,}2$ kbit/s útiles. Estos resultados excluyen
telemetría, sobrecarga del protocolo y retransmisiones. El tamaño real
del archivo deberá comprobarse con el formato de adquisición.


\subsection{Computador de a bordo e interfaz de alimentación}

Se adopta como referencia de diseño el Toradex Verdin iMX8M Plus Quad
4GB IT V1.1B, número de producto \texttt{00631101}. El fabricante
documenta esta referencia como reemplazo de la revisión V1.1A
\cite{verdinpcn}. Se elige la variante sin Wi-Fi/Bluetooth porque la
comunicación de misión corresponde al bloque de radio de la PCB.

El módulo integra 4 GB de RAM y 32 GB de eMMC; esta última se reserva
para arranque, sistema operativo y aplicación. Su puerto
\texttt{USB\_2} ofrece USB 3.1 Gen 1 en función anfitrión y se reserva
para la cámara. El nombre del puerto no indica velocidad USB 2.0.
\texttt{USB\_1} se reserva para recuperación. El rango industrial
es de $-40$ a $+85\,^{\circ}\mathrm{C}$, sujeto a las condiciones de
disipación \cite{verdinds}.

\subsubsection{Recursos y compatibilidad de software}
Dos buffers para cuadros RAW16 completos requieren
\begin{equation}
M_{\mathrm{buffers}}=2(2048\times1088\times2)
=8\,912\,896\ \mathrm{bytes}.
\end{equation}
Los 4 GB constituyen una configuración de catálogo con margen; no se
deducen como mínimo necesario de este cálculo. El uso total de RAM
incluye sistema operativo, SDK, copias y operaciones de procesamiento.

Se considera Linux ARM64 mediante un BSP de Toradex compatible con la
revisión V1.1B. Para la rama BSP 6, el aviso de cambio establece soporte
desde 6.2.0 \cite{verdinpcn}. Esta condición no demuestra por sí sola
compatibilidad con xiAPI: XIMEA ofrece software Linux ARM, pero no
enumera Verdin entre sus plataformas probadas \cite{ximealinux}.
El funcionamiento de esta combinación se mantiene como hipótesis
explícita de integración del diseño académico. No se ha ejecutado una
prueba ni verificado el conjunto de dependencias del SDK. La validación
posterior deberá fijar versiones y comprobar apertura de la cámara,
captura RAW16, integridad de archivos, tiempos y uso de memoria.

\subsubsection{Dimensionamiento de la rama de potencia}
La entrada recomendada del módulo admite de 3,135 a 5,5 V
\cite{verdinds}. Se adopta un riel de 5 V por coherencia con ET-02;
su intervalo de $\pm5\,\%$, de 4,75 a 5,25 V, está contenido en ese
rango. La tensión debe garantizarse en los pines VCC, considerando
las caídas de la distribución y protección.

La guía de familia establece una envolvente de 8,25 W sostenidos y
12,5 W de pico, y recomienda dimensionar la fuente para el pico
\cite{verdinguide}. Se utiliza esa envolvente como criterio de diseño,
sin atribuirla al consumo medido de la aplicación. En el extremo
inferior de tensión:
\begin{align}
I_{\mathrm{sost}}&=\frac{8{,}25}{4{,}75}
\approx1{,}737\ \mathrm{A},\\
I_{\mathrm{pico}}&=\frac{12{,}5}{4{,}75}
\approx2{,}632\ \mathrm{A}.
\end{align}
Se especifica una capacidad de 3 A para la rama del SoM. A 4,75 V
puede entregar 14,25 W, un margen del 14\,\% sobre la envolvente de
pico. La capacidad nominal de 15 W a 5 V no equivale al consumo
permanente del módulo. La cámara y los demás dispositivos requieren
presupuestos adicionales; la eMMC y RAM internas no se suman de nuevo.

La corriente especificada no sustituye la comprobación de rampa,
transitorios, desacoplo y secuencia de encendido. El diseño de la rama
debe satisfacer la gestión de potencia de la familia
\cite{verdinguide}. La energía se calculará a partir de consumos y
duraciones por modo, sin sustituirlos por la capacidad de la fuente.

\subsubsection{Integración en la PCB}
El módulo tiene una proyección de $69{,}6\times35$ mm
\cite{verdinds}. Su área proyectada es 2436 mm$^2$, equivalente al
30,07\,\% de una PCB de $90\times90$ mm. Esta comparación no incluye
zócalo, fijaciones, disipación ni separaciones, por lo que no demuestra
el encaje del conjunto completo.

Se reservarán circuitos para alimentación supervisada, recuperación,
USB de cámara y adaptación de niveles. Muchas señales digitales usan
1,8 V \cite{verdinguide}; cada interfaz deberá comprobarse antes de
conectarla a dispositivos de 3,3 V. La energía de la cámara se entrega
externamente; VBUS y la prevención de retroalimentación se cerrarán
según el contrato de la sección~\ref{sec:contrato_interfaces}.
El almacenamiento externo se conecta al bus SD dedicado. La interfaz
digital de comunicaciones usa SPI\_1; la asignación física debe comprobar
que no existan conflictos de multiplexación.

\subsection{Almacenamiento independiente de misión}

Se selecciona como referencia de diseño la microSD industrial Kingston
\texttt{SDCIT2/16GBSP}, de 16 GB nominales, instalada en un zócalo de la
PCB. La presentación SP no incluye adaptador SD. Su alimentación
nominal es de 3,3 V y su rango publicado de operación es de
$-40$ a $+85\,^{\circ}\mathrm{C}$ \cite{kingstonsdcit2}.
La memoria conserva imágenes, telemetría y resultados; la eMMC del
Verdin permanece dedicada al arranque, sistema operativo y aplicación.

Se prioriza microSD por la disponibilidad de una interfaz SD externa
en el computador y por permitir sustitución durante mantenimiento.
Una eMMC adicional soldada evitaría contactos extraíbles, pero requiere
su propia selección de encapsulado y circuito. La carga de datos
considerada no justifica incorporar un SSD con interfaz PCIe.
Esta elección exige asegurar la retención mecánica de la tarjeta y
no constituye una validación para operación orbital.

\subsubsection{Escenario de retención y justificación de capacidad}
Se establece una contingencia de 90 días sin descarga para dimensionar
la memoria. Es un supuesto de diseño, no un resultado orbital ni una
exigencia de la guía. Se conserva la adquisición de diez cuadros RAW16
diarios y se asignan las siguientes cuotas y reservas:
\begin{itemize}
\item Imágenes originales: $44{,}56448$ MB/día, según el cálculo de cámara.
\item Resultados de procesamiento: hasta un volumen adicional igual al
RAW. Es una cuota de almacenamiento; no define un algoritmo ni una
copia redundante de los originales.
\item Metadatos, telemetría, índices y eventos: cuota conjunta de
2 MB/día. No es una medición ni sustituye el dimensionamiento posterior
de registros y comunicaciones.
\item Ocupación inicial: máximo 1 GB al comenzar la contingencia.
La gestión de descarga y depuración deberá mantener esa condición;
si no se cumple, se recalculará la autonomía disponible.
\item Reserva libre: 3 GB para operación y contingencias, adicional
a los datos producidos durante el intervalo.
\end{itemize}
Se emplean unidades decimales: $1\ \mathrm{MB}=10^6$ bytes y
$1\ \mathrm{GB}=10^9$ bytes. La generación presupuestada es
\begin{equation}
\dot D_{\mathrm{pres}}=44{,}56448+44{,}56448+2
=91{,}12896\ \mathrm{MB/día}.
\end{equation}
Aplicando el criterio de retención definido previamente:
\begin{equation}
C_{\mathrm{req}}=1+\frac{91{,}12896(90)}{1000}+3
=12{,}2016064\ \mathrm{GB}.
\end{equation}
Se exige comprobar al menos 15 GB útiles después del formato y de
la configuración. Este valor es un umbral de aceptación, no una
capacidad medida. El fabricante advierte que parte de la capacidad
nominal no está disponible para datos \cite{kingstonsdcit2}.
Con dicho umbral, el margen adicional es
\begin{equation}
\Delta C=15-12{,}2016064=2{,}7983936\ \mathrm{GB}.
\end{equation}
Los 3 GB de reserva ya están incluidos en $C_{\mathrm{req}}$ y no se
suman nuevamente. La sensibilidad al período sin descarga es:
\begin{center}
\begin{tabular}{rrl}
\toprule
Retención (días) & Capacidad útil requerida (GB) & Frente a 15 GB útiles \\
\midrule
30 & 6,734 & Cumple el escenario \\
90 & 12,202 & Cumple el escenario \\
180 & 20,403 & No cumple \\
\bottomrule
\end{tabular}
\end{center}
Así se conserva ET-04, condicionado al cumplimiento de las cuotas y
del umbral útil. Si cambian los cuadros diarios, resultados, registros
o datos acumulados, se actualiza el cálculo. Más memoria no resuelve
una descarga permanentemente inferior a la generación.

\subsubsection{Interfaz y objetivo de escritura}
Se utiliza la interfaz externa \texttt{SD\_1}/USDHC2 del Verdin,
independiente de su eMMC interna. Se adopta inicialmente el modo
\emph{Default Speed}, cuatro bits a 25 MHz con señalización de 3,3 V
\cite{verdinds}. Esta elección limita las exigencias de velocidad y
evita depender de la conmutación UHS a 1,8 V. Toradex documenta el
acceso a tarjetas SD desde Linux en este módulo \cite{toradexsd}.

El techo bruto del bus es
\begin{equation}
R_{\mathrm{bus}}=\frac{4(25\times10^6)}{8}
=12{,}5\ \mathrm{MB/s}.
\end{equation}
No equivale a la velocidad sostenida del archivo. Como objetivo se
plantea escribir un RAW y un resultado de igual tamaño en hasta 2 s:
\begin{equation}
R_{\mathrm{objetivo}}=\frac{2(4\,456\,448)}{2}
=4{,}456448\ \mathrm{MB/s}.
\end{equation}
El cumplimiento deberá verificarse con el sistema de archivos y la
configuración elegidos, incluyendo bloqueos y ráfagas. No se infiere
del máximo comercial de escritura de la tarjeta.

\subsubsection{Alimentación e integridad de datos}
La rama de 3,3 V se dimensionará fuera de los 3 A reservados al SoM.
La ficha consultada no publica corriente máxima ni de arranque; por
ello la potencia máxima permanece sin valor, en vez de asumir cero.
El interruptor de alimentación deberá reconocer la habilitación
\texttt{SD\_1\_PWR\_EN}, que opera a 1,8 V aunque el bus SD use
3,3 V \cite{verdinds}. La elección del zócalo, desacoplo, protección
ESD e interruptor se completará con esos límites.

Se deberá evitar el corte de energía con escrituras activas, cerrar
y sincronizar archivos y comprobar su integridad. Las funciones de
protección publicadas por Kingston no acreditan por sí solas la
conservación del último archivo ni la consistencia del sistema de
archivos ante cualquier interrupción. Guardar RAW y resultados en
una misma tarjeta tampoco proporciona redundancia frente a su fallo.

La escritura lógica anual presupuestada es
$91{,}12896(365)/1000=33{,}2620704$ GB. No equivale a escritura física
de NAND; faltan amplificación y operaciones internas. No se atribuye
a la referencia de 16 GB el TBW máximo de la familia, que el fabricante
asocia a su capacidad mayor \cite{kingstonsdcit2}.

Las entradas, fórmulas, escenarios y restricciones correspondientes
se mantienen en el libro de cálculos, con su procedencia y condición
de dato publicado, supuesto o verificación pendiente.

\begin{thebibliography}{99}
\bibitem{guia4} Departamento de Ingeniería Electrónica, Universidad de Antioquia. \emph{Práctica 4: Diseño Electrónico y Captura Esquemática}. Archivo suministrado: \texttt{lab004\_2598575\_20262.pdf}. Véanse descripción, procedimiento 1 y entregables.
\bibitem{lab3} D. F. Yépez Taimal y S. Cáceres Vásquez. \emph{Práctica 3: Requisitos, especificaciones y arquitectura electrónica del CubeSat UD34}. 28 de septiembre de 2026. Archivo suministrado: \texttt{Diseño\_Lab3.pdf}. Véanse secciones 2--4 y 9--10; cuadros 16--19.

\bibitem{ximea7nir2} XIMEA. \emph{MQ022HG-IM-SM5X5-7NIR2: especificaciones del modelo}.
\url{https://www.ximea.com/products/hyperspectral-imaging/xispec-hyperspectral-miniature-cameras/imec-sm-range-600-1000-usb3-hyperspectral-camera}.
Consulta: 24 de octubre de 2026.

\bibitem{ximeamanual} XIMEA. \emph{xiSpec2 Technical Documentation}.
Versión 2.03 según portada, septiembre de 2026; pies de página V2.00.
Secciones 3.4, 4.1 y 10.1.
\url{https://www.ximea.com/getattachment/b750b3ca-2b30-42bf-bc83-4c5261ef52c6/xiSpec-Hyperspectral-HSI-mini-cameras_Technical-manual-DWL_manual.pdf}.
Consulta: 24 de octubre de 2026.

\bibitem{ximealinux} XIMEA. \emph{XIMEA Linux Software Package}, Supported platforms.
\url{https://www.ximea.com/support/wiki/apis/XIMEA_Linux_Software_Package}.
Consulta: 24 de octubre de 2026.

\bibitem{verdinds} Toradex. \emph{Verdin iMX8M Plus V1.1 Datasheet}.
Revisión documental 1.8, 22 de junio de 2026. Tablas 7, 30 y 72;
secciones 5.5 y 8.5--8.6.
\url{https://docs.toradex.com/116795-verdin_imx8mp_datasheet.pdf}.
Consulta: 24 de octubre de 2026.

\bibitem{verdinpcn} Toradex. \emph{PCN: 0063 Verdin iMX8M Plus Quad
4GB IT V1.1A to V1.1B}. Actualizado el 12 de diciembre de 2023.
\url{https://docs.toradex.com/114344-pcn-0063-v11a-v11b-20231127.pdf}.
Consulta: 24 de octubre de 2026.

\bibitem{verdinguide} Toradex. \emph{Verdin Carrier Board Design Guide}.
Versión 1.2, 15 de septiembre de 2022. Sección 3, especialmente 3.1.
\url{https://docs.toradex.com/108140-verdin-carrier-board-design-guide.pdf}.
Consulta: 24 de octubre de 2026.

\bibitem{kingstonsdcit2} Kingston. \emph{Industrial microSD Memory Card},
serie SDCIT2. Ficha de dos páginas; identificador MKD-02122026.
\url{https://www.kingston.com/datasheets/sdcit2_en.pdf}.
Consulta: 25 de octubre de 2026.

\bibitem{toradexsd} Toradex. \emph{SD/MMC Card (Linux)}.
\url{https://developer.toradex.com/linux-bsp/application-development/peripheral-access/sdmmc-card-linux/}.
Consulta: 25 de octubre de 2026.

\bibitem{jstgh}
J.S.T. Mfg. Co., Ltd., \emph{GH connector: Product Profile}.
\url{https://www.jst-mfg.com/product/index.php?lang=2&series=105}.
Consultado el 25 de octubre de 2026.
\bibitem{txu0202}
Texas Instruments, \emph{TXU0202}, SCES942A, revisión de marzo de 2022.
Secciones 6, 7.3, 7.5, 9.3 y 12.1.
\url{https://www.ti.com/lit/ds/symlink/txu0202.pdf}.
Consultado el 25 de octubre de 2026.

\bibitem{axc2t245}
Texas Instruments, \emph{SN74AXC2T245}, SCES879A, febrero de 2024.
\url{https://www.ti.com/lit/ds/symlink/sn74axc2t245.pdf}.
Consultado el 25 de octubre de 2026.
\bibitem{imx8mpiec} NXP Semiconductors, \emph{i.MX 8M Plus Applications
Processor Datasheet for Industrial Products}, IMX8MPIEC, rev. 2,
febrero de 2023, tabla 26, p. 38. Copia del documento del fabricante:
\url{https://static6.arrow.com/aropdfconversion/cb2b2904debebb012b3268b4bb38d817cef2cc4d/imx8mpiec.pdf}.
\bibitem{yageo10k} YAGEO, \emph{RC0603FR-0710KL, specification sheet}.
\url{https://www.yageogroup.com/component-documentation/download/specsheet/RC0603FR-0710KL}.
\bibitem{yageo47k} YAGEO, \emph{RC0603FR-0747KL, specification sheet}.
\url{https://www.yageogroup.com/component-documentation/download/specsheet/RC0603FR-0747KL}.
\bibitem{kemet100n} KEMET/YAGEO Group, \emph{C0603C104K3RACTU,
specification sheet}.
\url{https://yageogroup.com/component-documentation/download/specsheet/C0603C104K3RACTU?lang=en}.
\bibitem{esd321} Texas Instruments, \emph{ESD321}, SLVSEN8C, junio de
2025, secciones 4, 5.4--5.6 y 7.4.
\url{https://www.ti.com/lit/ds/symlink/esd321.pdf}.
Fuentes consultadas el 25 de octubre de 2026.
\bibitem{tps25947} Texas Instruments, \emph{TPS25947xx},
SLVSFC9C, revisión C, mayo de 2026, secciones 4 y 6.5.
\url{https://www.ti.com/lit/ds/symlink/tps25947.pdf}.
Consultado el 26 de octubre de 2026.
\bibitem{tps25948} Texas Instruments, \emph{TPS25948xx},
SLVSGT9D, revisión D, abril de 2026, secciones 4--6.5.
\url{https://www.ti.com/lit/ds/symlink/tps25948.pdf}.
Consultado el 26 de octubre de 2026.
\end{thebibliography}
\end{document}
